\documentclass[twocolumn]{aastex631}
\begin{document}
\title{The reionization ionizing-photon-budget shortfall for star-forming galaxies is not robust to systematics at z\~{}7 (jwsthigh systematic corner)}
\author{NebulaMind Lab (autonomous pipeline)}
\affiliation{NebulaMind Lab, \url{https://lab.nebulamind.net}}
\begin{abstract}
Reionization ionizing-photon-budget reconciliation at z\~{}7: star-forming galaxies require f\_esc=0.040 (+0.038/-0.019) to close the budget (Madau-Dickinson SFRD, log xi\_ion=25.5+/-0.15, clumping C=2-5, JWST-SFRD tail), versus indirect-proxy-inferred f\_esc=0.062 (+0.110/-0.039) from LzLCS O32/beta calibrations. Median delta(required-inferred)=-0.019 dex-frac (16-84\%: -0.128 to +0.029); 36\% of the systematic MC shows a shortfall. Conclusion: the budget CLOSES within the systematic, and the sign holds under both O32 and beta calibrations. The result is bounded by the xi\_ion x clumping x proxy-calibration systematic, not by statistics. Generated autonomously from public data (jwst) via the NebulaMind Lab runner: a bounded, reproducible, descriptive study.
\end{abstract}
\section{Introduction}
Introduction
The reionization of the universe remains a topic of significant interest in astrophysics. Recent studies have highlighted potential discrepancies between the observed star formation rate density (SFRD) and the required ionizing photon budget to achieve reionization [Muoz2024, Davies2021]. This has led to questions about whether star-forming galaxies alone can account for the necessary ionizing photons. To address this issue, we revisit the reionization-photon-budget using a literature-anchored approach.
\section{Data and method}
Data and method
Our analysis relies on published values of cosmic SFRD [Madau2017], ionizing efficiency (xi\_ion), and escape fraction (f\_esc) calibrations from Lyman-alpha-emitting galaxies at high redshifts [LzLCS: Chisholm+22, Flury+22; Simmonds+24]. We adopt the Madau \& Dickinson (2014) analytic fitting function for SFRD and use the O32/beta f\_esc proxy calibrations to estimate the ionizing photon budget. Our method focuses on reconciling the systematic uncertainties in these parameters to determine if star-forming galaxies can close the reionization-photon-budget at z\~{}7.
\section{Result}
Result
Reconciling the reionization ionizing-photon-budget at z\~{}7, we find that star-forming galaxies require an escape fraction of f\_esc=0.040 (+0.038/-0.019) to match the Madau-Dickinson SFRD with log xi\_ion=25.5±0.15 and clumping factor C=2-5. This is compared to the indirect-proxy-inferred f\_esc=0.062 (+0.110/-0.039) from LzLCS O32/beta calibrations. The median difference between required and inferred escape fractions is -0.019 dex-frac, with a 16-84\% range of -0.128 to +0.029. Notably, 36\% of our systematic Monte Carlo simulations show a shortfall in the ionizing photon budget.
\begin{figure}\centering\includegraphics[width=\columnwidth]{result.png}\caption{The reionization ionizing-photon-budget shortfall for star-forming galaxies is not robust to systematics at z\~{}7 (jwsthigh systematic corner)}\end{figure}
\section{Caveats}
Caveats
Our study relies on an automated, single-selection, uncalibrated measurement approach, which has inherent limitations. The accuracy of our results depends heavily on the adopted literature values and calibrations, which may not fully capture the complexities of reionization processes. Furthermore, uncertainties in xi\_ion, clumping factor C, and proxy calibrations can significantly impact the estimated escape fraction. Additionally, our analysis does not account for potential contributions from other ionizing sources, such as active galactic nuclei or quasars. Therefore, while our results provide valuable insights into the reionization-photon-budget, they should be interpreted with caution and considered alongside other observational and theoretical studies to obtain a comprehensive understanding of this complex phenomenon.
\section*{References}
\footnotesize
[Muoz2024] Muñoz et al. (2024). Reionization after JWST: a photon budget crisis?. bibcode 2024MNRAS.535L..37M \par [Park2022] Park et al. (2022). Calibrating excursion set reionization models to approximately conserve ionizing photons. bibcode 2022MNRAS.517..192P \par [Duncan2015] Duncan et al. (2015). Powering reionization: assessing the galaxy ionizing photon budget at z < 10. bibcode 2015MNRAS.451.2030D \par [Davies2021] Davies et al. (2021). The Predicament of Absorption-dominated Reionization: Increased Demands on Ionizing Sources. bibcode 2021ApJ...918L..35D \par [Madau2017] Madau (2017). Cosmic Reionization after Planck and before JWST: An Analytic Approach. bibcode 2017ApJ...851...50M

\end{document}
